% Gerado por tcc/gerar_latex.py a partir de TCC_exo1_exo2.docx.
% Edite à vontade; uma nova geração sobrescreve este arquivo.

\chapter{Conclusão}\label{sec:5}

As conclusões estão organizadas por dispositivo: a seção~\ref{sec:5.1} trata do gerador de torque da US 9.889.554 B2, e a seção~\ref{sec:5.2}, do dispositivo da EP 3 504 033 B1.

\section{Exoesqueleto 1}\label{sec:5.1}

Este trabalho reconstruiu, a partir da documentação de patente, o gerador de torque de mola da US 9.889.554 B2 e formulou um modelo estático do torque que ele entrega em função do ângulo. A cinemática foi resolvida em forma fechada a partir da equação de fechamento do circuito vetorial, e cada etapa da cadeia de cálculo foi conferida por verificações independentes, com resíduos da ordem de 10\textsuperscript{−13} mm nas relações geométricas.

O princípio dos trabalhos virtuais mostrou-se o critério decisivo na formulação. Ele revelou que a extensão da mola deve acompanhar o caminho completo do cabo, ℓ = s − Rψ, e não apenas o vão livre, que violava a conservação de energia com resíduo de 1,71 mm. A mesma verificação distinguiu o braço de momento perpendicular, limitado a 22,919 mm, da distância do pivô ao ponto de tangência, próxima de 283 mm, cuja confusão superestimaria o torque por um fator superior a 12.

Com os parâmetros congelados, o modelo reproduz a forma da curva publicada com coeficiente de determinação de 0,99839 e RMSE de 1,83 u.p., equivalente a 1,24\% do maior valor extraído. Os resíduos, contudo, organizam-se em zonas de sinal alternado, o que indica diferença sistemática de forma, e a formulação anterior, fisicamente inconsistente, ajustava melhor a curva. A adoção da formulação consistente, apesar do RMSE maior, reflete o critério de que a escolha entre modelos se faz por validade física, e não pelo indicador de erro.

A varredura de h com K congelado separou o que a geometria impõe do que a curva extraída sugere. O cabo permanece tracionado em toda a faixa de operação somente para h ≤ 32,886 mm, o que exclui a cota de 34,785 mm lida no CAD. A curva extraída favorece um intervalo de 0,145 mm em torno do valor adotado, mas essa concentração é esperada por construção, porque K e X\textsubscript{pose} foram identificados nesse mesmo valor; o critério que levou a h = 22,701 mm não pôde ser reproduzido, e a afirmação anterior de decrescimento monotônico do erro com h foi refutada pela auditoria.

No confronto com o torque gravitacional do conjunto braço–exoesqueleto, cujo pico de 11,16 N·m ocorre a 89,4° de elevação, as curvas normalizadas diferem em no máximo 0,027 entre 105° e 150° de elevação. O pico do torque do cabo, a 97,5°, fica 8,05° acima do pico gravitacional, e o ponto morto, a 30,5°, libera o braço nas posturas de não trabalho, à custa de subcompensação abaixo de 94°. O resultado quantifica o compromisso de um mecanismo passivo de torque fixo por ângulo, que favorece o trabalho acima da cabeça.

O objetivo geral foi atingido em sua parte estática e para o primeiro dos três dispositivos previstos; o segundo é tratado na seção~\ref{sec:5.2}. Os objetivos específicos a, b, c e e foram cumpridos; o objetivo d, parcialmente, pois a rigidez e a extensão da mola na pose foram reproduzidas, mas a origem da posição da ancoragem e da calibração angular não. Os objetivos f e g, referentes à equação de movimento e ao acoplamento com o MBDyn, não foram executados. A principal restrição é a ausência de unidade na ordenada da curva da patente: o torque do cabo permanece em unidade da patente, e a leitura em N·m conduz a forças implausíveis, de modo que o modelo determina a forma do torque, mas não sua escala física.

Como continuidade, recomenda-se obter um dado independente que fixe a escala do torque, como a especificação de uma mola comercial compatível; conferir no CAD a cota de ancoragem, o comprimento do alojamento e o comprimento da mola na pose; documentar o procedimento de extração dos pontos e sua incerteza; representar a protuberância responsável pela zona neutra; e, com o torque em unidade física, executar as etapas dinâmicas pela formulação de Lagrange e no MBDyn e estender o procedimento à US 9.737.374 B2. A extensão à EP 3 504 033 B1 é apresentada nas seções~\ref{sec:3.2}, \ref{sec:4.2} e \ref{sec:5.2}.

\section{Exoesqueleto 2}\label{sec:5.2}

Para a EP 3 504 033 B1, foram executadas a leitura da arquitetura na patente, a organização das grandezas do CAD com sua origem, a composição inercial do membro humano, as convenções angulares, a cinemática de dois regimes em forma fechada, o torque pelo produto vetorial, o torque gravitacional, o tratamento da Figura 16 da patente, a identificação de s\textsubscript{0} pelo espaçamento entre as curvas e pelo ajuste de forma, a varredura do trecho e\textsubscript{c} e as verificações numéricas.

O espaçamento entre as três curvas da Figura 16 da patente determina os comprimentos de referência do elástico sem depender da escala da ordenada, da rigidez, do braço de momento ou do torque do peso: 300,91 mm na curva pontilhada, 287,19 mm na tracejada e 275,59 mm na contínua. Na curva tracejada, que fornece o s\textsubscript{0} adotado, a variação entre as faixas de ângulo testadas é de até 2,73 mm, e o desvio padrão com o ruído de extração, de 1,02 mm (0,43 mm na pontilhada e 1,43 mm na contínua). O fim da curva pontilhada, de pré-tensão quase nula, confirma o resultado de forma independente: a razão dessa curva, extrapolada, se anula em 171,2°, o que o modelo interpreta como afrouxamento do elástico, e com o s\textsubscript{0} identificado o modelo afrouxa a partir de 166,4°. O início das curvas, perto de 14°, não informa s\textsubscript{0} e foi atribuído ao contato do elástico com o elemento de encosto junto à articulação.

A forma do torque prevista com a geometria do CAD e o peso do membro humano não reproduz as curvas da patente: o fator comum às três curvas difere do previsto por um fator que vai de 0,39 em 20° a 2,97 em 165°, normalizado em 90°. O ajuste de forma, que deixa s\textsubscript{0} e a escala livres por curva, compensa essa diferença com valores de s\textsubscript{0} sem significado físico, como o limite s\textsubscript{0} = 0 na curva contínua. A origem da diferença não pôde ser determinada; entre 35° e 155°, os dados favorecem o torque do peso usado pela patente para formar a razão, ou a definição da ordenada, sobre o braço de momento.

A rigidez não é identificável pela figura. Sob a hipótese de assistência de 40\% do torque do peso em 90°, K = 1,467 N/mm, e o elástico trabalha de 149,7 N com o braço pendente a 14,3 N com o braço elevado, com deformação máxima de 35,5\% do comprimento de referência. O torque atinge 4,376 N·m em 78,8°, dentro da faixa de 70° a 120° em que a patente prefere a força máxima, lida aqui como força de apoio proporcional ao torque, e a razão τ/τ\textsubscript{g} vale 0,47 em 10°, 0,34 em 30°, 0,43 em 60°, 0,40 em 90° e 0,21 em 150°, com mínimo de 0,18 em 163,4°. Abaixo da transição de regime, em 19,09°, o modelo ainda entrega torque, porque o elemento de encosto do CAD está a 7,754 mm do eixo.

As relações de circuito dos dois regimes e o princípio dos trabalhos virtuais foram verificados com resíduos de até 1,1 × 10\textsuperscript{−13} mm e 7,6 × 10\textsuperscript{−8} mm/rad. Com s\textsubscript{0} e K mantidos, a varredura do trecho e\textsubscript{c} mostrou que o elástico permanece tracionado em toda a faixa só para e\textsubscript{c} ≤ 44,50 mm; o aumento de e\textsubscript{c} antecipa a transição de regime, de 26,91° para 14,46°, e desloca o pico do torque de 84,7° para 74,8° entre e\textsubscript{c} = 25 mm e e\textsubscript{c} = 45 mm. A comparação com a Figura 17 da patente ficou restrita à inspeção visual.

Dos objetivos acrescentados para esse dispositivo, o item h foi cumprido na parte estática, com três ressalvas: a escala do torque depende da hipótese K = K\textsubscript{40}; as grandezas do CAD não puderam ser conferidas; e a composição inercial e o torque gravitacional incluem o membro humano e o braço do exoesqueleto, mas não a alavanca, o elemento 78 e a braçadeira, previstos nos itens a e e. O item i foi cumprido pelo espaçamento entre as curvas. Para esse dispositivo, o objetivo geral foi atingido na parte estática, com o confronto com o torque gravitacional feito pela forma normalizada e, em N·m, apenas no cenário de projeto. As etapas dinâmicas não foram executadas, pelas razões da seção~\ref{sec:3.2.12}.

Os dois dispositivos diferem tanto na arquitetura quanto no que suas patentes permitem identificar. O gerador da US 9.889.554 B2 reproduz a curva publicada com R\textsuperscript{2} de 0,99839, mas em unidade da patente; o dispositivo da EP 3 504 033 B1 tem s\textsubscript{0} determinado de forma mais robusta, por dispor de três curvas, mas sua forma não é reproduzida pelo modelo do CAD, e o torque só chega a N·m por hipótese. Nos dois casos, o pico de assistência fica afastado do pico gravitacional, acima dele no primeiro (97,5° contra 89,4°) e abaixo no segundo (78,8° contra 89,7°), e a faixa sem torque com o braço baixado é produzida por meios distintos, uma protuberância no primeiro e o contato com o elemento de encosto no segundo. A comparação proposta na seção~\ref{sec:2.11} fica, assim, respondida na parte estática quanto aos ângulos e à forma dos perfis, mas não quanto à magnitude.

Como continuidade para esse dispositivo, recomenda-se incluir o modelo em CAD no repositório e confirmar nele o sinal do ângulo do braço na pose e o valor de e\textsubscript{c}, que coincide com a cota H\textsubscript{CAD} do outro mecanismo; esclarecer a definição da ordenada da Figura 16 da patente e a configuração do dispositivo que a gerou, que a patente não informa; conferir a posição e o raio de contato do elemento de encosto, comparando com a posição sobre o eixo preferida pela patente; registrar o procedimento e a incerteza da extração da Figura 16 da patente; obter a curva de força e extensão do elemento elástico, que fixaria K; digitalizar a Figura 17 da patente; e, com a escala resolvida, executar as etapas dinâmicas pela formulação de Lagrange e no MBDyn.
